\subsection{Persons (5123034)}
Die Schnittstelle für den Bereich \textit{Persons} im Backend umfasst die Verwaltung von Patienten, Ärzten sowie Pflegepersonal. Dabei bildet eine zentrale Personentabelle die Grundlage, auf der die verschiedenen Rollen (Patient, Mitarbeiter) modelliert werden.
Ein Großteil der Funktionalität in diesem Modul besteht aus der Bereitstellung von CRUD-Schnittstellen, wobei die Daten in vielen Fällen weitgehend unverändert zwischen Frontend und Datenbank weitergereicht werden. Dennoch existieren gezielte Erweiterungen, um Datenkonsistenz und Integrität sicherzustellen.
Zu den grundlegenden Funktionen zählen Suchschnittstellen, über die Personen anhand spezifischer Identifikatoren (z.,B. Patienten-ID, Doktor-ID oder Pfleger-ID) sowie anhand allgemeiner Attribute abgefragt werden können.
Darüber hinaus werden im Patienten-Kontext zusätzliche Funktionalitäten bereitgestellt. Hierzu gehören insbesondere die Verwaltung von Aufenthalten (Bookings), einschließlich der Anpassung und vorzeitigen Beendigung, sowie der Zugriff auf diagnostische Informationen.
\subsubsection{Methodik der Implementierung}
Die Implementierung folgt dem Ansatz einer zentralen Personenverwaltung mit rollenbasierter Erweiterung. Das bedeutet, dass jede Person zunächst unabhängig von ihrer Rolle in einer gemeinsamen Tabelle gespeichert wird und anschließend über zusätzliche Entitäten (z.,B. Patient oder Employee) spezifiziert wird.
Eine besondere Herausforderung stellt das Anlegen neuer Personen dar. Um doppelte Einträge in der zentralen Personentabelle zu vermeiden, wurde ein mehrstufiger Prüfmechanismus implementiert:
Bei einem POST-Request ohne übergebene UUID wird zunächst eine Ähnlichkeitssuche auf Basis ausgewählter Attribute (z.,B. Name, Geburtsdatum) durchgeführt, um potenzielle bereits existierende Personen zu identifizieren. Werden mögliche Treffer gefunden, so werden diese mit relevanten Unterscheidungsmerkmalen an das Frontend zurückgegeben, sodass eine bewusste Auswahl getroffen werden kann.
Erst wenn keine passenden Einträge existieren oder explizit eine neue Person angelegt werden soll, erfolgt die Erstellung eines neuen Datensatzes inklusive Generierung einer eindeutigen UUID.
Zusätzlich wird beim Anlegen spezifischer Rollen (z.,B. Arzt) sichergestellt, dass alle abhängigen Entitäten konsistent erstellt werden. Beispielsweise wird beim Erstellen eines neuen Arztes automatisch auch ein entsprechender Employee-Eintrag generiert.
\subsubsection{Begründung der Vorgehensweise}
Die Verwendung einer zentralen Personentabelle mit rollenbasierter Erweiterung ermöglicht eine flexible und redundanzarme Datenhaltung. Insbesondere wird dadurch vermieden, dass identische Personen mehrfach in unterschiedlichen Rollen separat gespeichert werden.
Der implementierte Mechanismus zur Duplikatserkennung dient der Sicherstellung einer hohen Datenqualität. In realen Anwendungsszenarien ist es häufig möglich, dass eine Person bereits im System existiert (z.,B. ein Mitarbeiter, der später als Patient aufgenommen wird). Ohne entsprechende Prüfung würde dies zu Inkonsistenzen und redundanten Daten führen.
Die Entscheidung, potenzielle Duplikate zunächst an das Frontend zurückzumelden, erlaubt eine kontrollierte Interaktion mit dem Nutzer und verhindert fehlerhafte automatische Zusammenführungen.
Die automatische Erstellung abhängiger Entitäten (z.,B. Employee beim Anlegen eines Arztes) reduziert zudem die Fehleranfälligkeit der API-Nutzung und stellt sicher, dass das Datenmodell jederzeit in einem konsistenten Zustand bleibt.
Insgesamt trägt dieser Ansatz dazu bei, die Datenintegrität zu gewährleisten und gleichzeitig eine flexible Erweiterbarkeit des Systems zu ermöglichen.
